\chapter{Training}

\section{Pipeline di addestramento}

Il training loop \`e implementato nel modulo \texttt{trainer.py} e supporta tutte le architetture tramite un'interfaccia unificata. La funzione \texttt{run\_training} orchestra l'intero processo: warmup del learning rate, alternanza train/validation, checkpointing, early stopping e logging.

\subsection{Mixed Precision (AMP)}

Tutti gli esperimenti (eccetto il Denoising Autoencoder) utilizzano Automatic Mixed Precision di PyTorch: le operazioni nel forward pass vengono eseguite in FP16 con \texttt{torch.autocast}, mentre la loss viene sempre calcolata in FP32 per stabilit\`a numerica. Un \texttt{GradScaler} gestisce il ridimensionamento dei gradienti.

La validazione viene eseguita \textbf{sempre in full precision}, indipendentemente dall'impostazione AMP, per evitare instabilit\`a causate dall'interazione tra BatchNorm in modalit\`a eval e FP16.

\subsection{Learning Rate Warmup}

I primi \texttt{warmup\_epochs} (tipicamente 5--8) utilizzano un warmup lineare del learning rate da \texttt{warmup\_start\_lr} (o $\text{lr} \times \frac{\text{epoch}}{\text{warmup\_epochs}}$ se non specificato) fino al learning rate target. Lo scheduler principale (CosineAnnealing o StepLR) viene attivato solo al termine del warmup.

\subsection{Gradient Clipping}

I gradienti vengono clippati tramite \texttt{clip\_grad\_norm\_} con norma massima configurabile (tipicamente 0.5--1.0). In caso di AMP, i gradienti vengono prima ``unscaled'' dal \texttt{GradScaler} prima del clipping.

\subsection{Early Stopping}

Il contatore di patience viene incrementato solo nelle epoche di validazione in cui la val loss non migliora. Questo permette di configurare la validazione a intervalli (es.\ ogni 2 epoche) senza penalizzare il conteggio.

\subsection{Gestione errori OOM}

In caso di errore \textit{Out of Memory} durante il training, il sistema:
\begin{enumerate}
    \item Salva un checkpoint di emergenza
    \item Libera la memoria CUDA
    \item Stampa un messaggio diagnostico con suggerimenti (ridurre batch/patch size, abilitare AMP)
\end{enumerate}

% ============================================================
\section{Training Pix2Pix}

Il training del Pix2Pix segue il paradigma GAN con addestramento alternato di generatore e discriminatore.

\subsection{Training del Discriminatore}

Il discriminatore riceve coppie di immagini (condizionale + target):
\begin{itemize}
    \item \textbf{Coppie reali:} $(x_{\text{degraded}},\; x_{\text{clean}}) \to$ label 1
    \item \textbf{Coppie false:} $(x_{\text{degraded}},\; G(x_{\text{degraded}})) \to$ label 0
\end{itemize}

La loss del discriminatore \`e:
$$
\mathcal{L}_D = \frac{1}{2}\left(\text{BCE}(D(x, y_{\text{real}}), 1) + \text{BCE}(D(x, y_{\text{fake}}), 0)\right)
$$

\subsection{Training del Generatore}

La loss del generatore combina la componente avversariale con una loss di ricostruzione:
$$
\mathcal{L}_G = \text{BCE}(D(x, G(x)), 1) + \lambda_{\text{L1}} \cdot \mathcal{L}_{\text{reconstruction}}(G(x), y)
$$

dove $\lambda_{\text{L1}} = 150$ e la loss di ricostruzione \`e la \texttt{CombinedLoss} (L1 + SSIM). Il fattore $\lambda_{\text{L1}}$ elevato bilancia la componente avversariale, che tende a dominare nelle prime fasi del training.

\subsection{Diagnostica del rumore}

In modalit\`a residual learning, vengono monitorate statistiche sul rumore predetto: media, deviazione standard, valore massimo, percentuale di outlier (valori che superano il clamp prima del clamping) e percentuale di pixel clippati nell'output.

% ============================================================
\section{Configurazioni degli esperimenti}

La tabella seguente riassume gli iperparametri principali di ciascun esperimento condotto. Tutti gli esperimenti utilizzano il dataset DIV2K con rumore gaussiano $\sigma = 100$, salvo diversa indicazione.

\begin{table}[H]
\centering
\small
\setlength{\tabcolsep}{4pt}
\begin{tabular}{lccccccc}
\toprule
\textbf{Modello} & \textbf{Loss} & \textbf{LR} & \textbf{Epoche} & \textbf{Batch} & \textbf{Patch} & \textbf{Warmup} & \textbf{Patience} \\
\midrule
UNet (bilinear) & L1+SSIM & 1e-4 & 36 & 16 & 128 & 5 & 5 \\
UNet Res. (bilinear) & L1+SSIM+VGG & 1e-4 & 36 & 16 & 128 & 5 & 5 \\
UNet Res. (transconv) & L1+SSIM+VGG & 1e-4 & 36 & 16 & 128 & 5 & 5 \\
Att.\ UNet (bilinear, all) & L1+SSIM+VGG & 4e-5 & 36 & 16 & 128 & 8 & 5 \\
Att.\ UNet (bilinear, btlnk) & L1+SSIM+VGG & 4e-5 & 36 & 16 & 128 & 8 & 3 \\
Att.\ UNet (transconv, all) & L1+SSIM+VGG & 4e-5 & 36 & 16 & 128 & 8 & 5 \\
Att.\ UNet (transconv, btlnk) & L1+SSIM+VGG & 4e-5 & 36 & 16 & 128 & 8 & 5 \\
\midrule
DnCNN-20 & L1+SSIM & 1e-4 & 100 & 64 & 80 & --- & 10 \\
Denoising AE & L1+SSIM & 1e-4 & 200 & 40 & 128 & 2 & 15 \\
RIDNet (8 EAB) & L1+SSIM & 1e-4 & 150 & 8 & 256 & 5 & 10 \\
\midrule
Pix2Pix & GAN+L1+SSIM & G:2e-4 & 150 & 8 & 384 & --- & 25 \\
\bottomrule
\end{tabular}
\caption{Iperparametri di training per ciascun esperimento. Le loss L1+SSIM usano pesi $\alpha\!=\!0.84, \beta\!=\!0.16$; le loss L1+SSIM+VGG usano $\alpha\!=\!0.6, \beta\!=\!0.25, \gamma\!=\!0.15$ (salvo varianti bottleneck con $0.65/0.3/0.05$).}
\label{tab:iperparametri}
\end{table}

\subsection{Note sugli esperimenti}

\begin{itemize}
    \item \textbf{Ottimizzatore:} Tutti i modelli usano AdamW ($\text{weight\_decay} = 10^{-5}$), eccetto DnCNN e RIDNet che usano Adam ($\text{weight\_decay} = 10^{-4}$) e Pix2Pix che usa Adam con $\beta_1 = 0.5$.
    
    \item \textbf{Scheduler:} CosineAnnealingLR per la maggior parte dei modelli; StepLR (step=1, $\gamma\!=\!0.96$) per DnCNN; ReduceLROnPlateau per il discriminatore Pix2Pix; nessuno scheduler per il generatore Pix2Pix.
    
    \item \textbf{Attention UNet --- varianti bottleneck:} Utilizzano GroupNorm invece di BatchNorm, layer percettuali meno profondi (\texttt{relu1\_2} + \texttt{relu2\_2} anzich\'e \texttt{relu2\_2} + \texttt{relu3\_3}) e peso percettuale ridotto ($\gamma = 0.05$).
    
    \item \textbf{RIDNet:} Utilizza 8 EAB con \texttt{residual\_scale} = 0.1, feature a 96 canali (anzich\'e 64) e gradient clip a 0.5 (pi\`u restrittivo della norma 1.0 usata dagli altri).
    
    \item \textbf{Dithering:} Sono stati condotti due esperimenti anche con degradazione di tipo quantizzazione a 2 bit con random dithering, usando la UNet base.
\end{itemize}
